\documentclass[11pt]{article}

\usepackage{fullpage}
\usepackage{amsmath}
\usepackage{amssymb}
\usepackage{amsfonts}
\usepackage{mathtools}
\usepackage[colorlinks,linkcolor=magenta,citecolor=blue]{hyperref}

\newcommand{\Rbar}{\bar R}
\newcommand{\Om}{\Omega}
\newcommand{\E}{\mathbb E}
\newcommand{\tr}{\operatorname{tr}}
\newcommand{\diag}{\operatorname{diag}}
\newcommand{\Unif}{\operatorname{Unif}}
\newcommand{\Normal}{\mathcal N}

\title{A Self-Contained Derivation of SER-relax}
\author{}
\date{}

\begin{document}
\maketitle

\section{Model}

Let \(x\in\mathbb R^J\) be the vector of marginal association
statistics and let \(N\) be the GWAS sample size. Let
\(\Rbar\in\mathbb R^{J\times J}\) be a positive definite reference LD
correlation matrix, and write
\[
\Om=\Rbar^{-1}.
\]
For each candidate active coordinate \(j\), define the reference
conditional mean and Schur complement
\[
\mu_j=\Rbar_{-j,j},
\qquad
S_j=\Rbar_{-j,-j}-\Rbar_{-j,j}\Rbar_{j,-j}.
\]
Because \(\Rbar\) is a correlation matrix, \(\Rbar_{jj}=1\). For
\(z_j\in\mathbb R^{J-1}\), let \(u_j(z_j)\in\mathbb R^J\) denote the
vector whose \(j\)-th entry is \(1\) and whose remaining entries are
\(z_j\).

SER-relax is a single-effect model that keeps the observation covariance
fixed at \(\Rbar/N\), but allows the active LD column to move around the
reference LD column. The model is
\begin{align}
\gamma &\sim \Unif\{1,\ldots,J\},\\
\beta &\sim \Normal(0,\sigma_0^2),\\
z_j\mid \gamma=j,\nu_0
&\sim \Normal_{J-1}(\mu_j,\tau^2S_j),\\
x\mid \beta,\gamma=j,z_j
&\sim \Normal_J\left(\beta u_j(z_j),\frac{\Rbar}{N}\right),
\end{align}
where
\[
\tau^2=\frac{1}{\nu_0+3},
\qquad
\tau^{-2}=\nu_0+3.
\]
Large \(\nu_0\) makes \(z_j\) concentrate near the reference column
\(\mu_j\). Small \(\nu_0\) allows more LD-column relaxation.

\section{Useful Block Identities}

The block inverse formula for \(\Rbar\) gives
\[
\Om_{-j,-j}=S_j^{-1},
\qquad
\Om_{-j,j}=-S_j^{-1}\mu_j.
\]
Therefore, for any \(z_j\), with
\[
\delta_j=x_{-j}-\mu_jx_j,
\qquad
y_j=z_j-\mu_j,
\]
we have
\begin{align}
u_j(z_j)^\top\Om u_j(z_j)
&=1+y_j^\top S_j^{-1}y_j, \label{eq:u-Omega-u}\\
u_j(z_j)^\top\Om x
&=x_j+y_j^\top S_j^{-1}\delta_j. \label{eq:u-Omega-x}
\end{align}
These two identities are the algebraic core of the CAVI updates.

\section{Variational Family}

Use the mean-field approximation
\[
q(\gamma=j,\beta,z_j)=\pi_jq_j(\beta)q_j(z_j),
\]
where
\[
q_j(\beta)=\Normal(m_{\beta j},v_{\beta j}),
\qquad
q_j(z_j)=\Normal_{J-1}(m_{zj},V_{zj}).
\]
Write
\[
\bar\beta_j=m_{\beta j},
\qquad
\overline{\beta^2}_j=v_{\beta j}+m_{\beta j}^2.
\]
For the current \(q_j(z_j)\), define
\begin{align}
Q_j
&=
\tr(S_j^{-1}V_{zj})
+(m_{zj}-\mu_j)^\top S_j^{-1}(m_{zj}-\mu_j),\\
A_j&=1+Q_j,\\
B_j
&=
x_j+(m_{zj}-\mu_j)^\top S_j^{-1}\delta_j.
\end{align}
Here \(Q_j=\E_j[y_j^\top S_j^{-1}y_j]\),
\(A_j=\E_j[u_j(z_j)^\top\Om u_j(z_j)]\), and
\(B_j=\E_j[u_j(z_j)^\top\Om x]\).

\section{CAVI Derivation}

Conditional on \(\gamma=j\), the log joint density, up to constants that
do not depend on \(\beta\) or \(z_j\), is
\begin{align}
\log p(x,\beta,z_j,\gamma=j)
&=
\log(1/J)
-\frac{N}{2}\{x-\beta u_j(z_j)\}^\top
\Om
\{x-\beta u_j(z_j)\} \nonumber\\
&\quad
-\frac{\beta^2}{2\sigma_0^2}
-\frac{\tau^{-2}}{2}(z_j-\mu_j)^\top S_j^{-1}(z_j-\mu_j).
\label{eq:log-joint}
\end{align}

\subsection{Update for \texorpdfstring{\(q_j(\beta)\)}{qj(beta)}}

The optimal coordinate update satisfies
\[
\log q_j^*(\beta)
=
\E_{q_j(z_j)}\log p(x,\beta,z_j,\gamma=j)+\mathrm{const}.
\]
Keeping only the terms that depend on \(\beta\), using
\eqref{eq:u-Omega-u}--\eqref{eq:u-Omega-x},
\begin{align}
\log q_j^*(\beta)
&=
N\beta B_j
-\frac12\left(\sigma_0^{-2}+NA_j\right)\beta^2
+\mathrm{const}.
\end{align}
Thus the update is Gaussian:
\begin{align}
v_{\beta j}
&=
\left(\sigma_0^{-2}+NA_j\right)^{-1},\\
m_{\beta j}
&=
v_{\beta j}NB_j.
\end{align}

\subsection{Update for \texorpdfstring{\(q_j(z_j)\)}{qj(zj)}}

The optimal relaxed-column update satisfies
\[
\log q_j^*(z_j)
=
\E_{q_j(\beta)}\log p(x,\beta,z_j,\gamma=j)+\mathrm{const}.
\]
Let \(y_j=z_j-\mu_j\). The terms depending on \(y_j\) are
\begin{align}
\log q_j^*(z_j)
&=
N\bar\beta_j y_j^\top S_j^{-1}\delta_j
-\frac12
\left(N\overline{\beta^2}_j+\tau^{-2}\right)
y_j^\top S_j^{-1}y_j
+\mathrm{const}.
\end{align}
Completing the square gives
\begin{align}
V_{zj}
&=
\frac{S_j}{N\overline{\beta^2}_j+\tau^{-2}},\\
m_{zj}
&=
\mu_j+
\frac{N\bar\beta_j}
{N\overline{\beta^2}_j+\tau^{-2}}
\delta_j.
\end{align}

\subsection{Update for \texorpdfstring{\(\pi_j\)}{pi j}}

The optimal categorical probabilities are proportional to the local
variational lower bounds:
\[
\pi_j
=
\frac{\exp(\mathcal L_j)}
{\sum_{k=1}^J\exp(\mathcal L_k)}.
\]
For SER-relax,
\begin{align}
\mathcal L_j
&=
\log(1/J)
+\E_j\log p(x\mid\beta,z_j,\gamma=j)
+\E_j\log p(\beta)
+\E_j\log p(z_j\mid\nu_0) \nonumber\\
&\quad
-\E_j\log q_j(\beta)
-\E_j\log q_j(z_j).
\end{align}
The expected log likelihood is
\begin{align}
\E_j\log p(x\mid\beta,z_j,\gamma=j)
&=
-\frac{J}{2}\log(2\pi)
-\frac12\log\left|\frac{\Rbar}{N}\right| \nonumber\\
&\quad
-\frac{N}{2}
\left[
x^\top\Om x
-2\bar\beta_jB_j
+\overline{\beta^2}_jA_j
\right].
\end{align}
The Gaussian prior and entropy terms are
\begin{align}
\E_j\log p(\beta)
&=
-\frac12\log(2\pi\sigma_0^2)
-\frac{\overline{\beta^2}_j}{2\sigma_0^2},\\
\E_j\log p(z_j\mid\nu_0)
&=
-\frac{J-1}{2}\log(2\pi)
-\frac12\log|S_j|
+\frac{J-1}{2}\log\tau^{-2}
-\frac{\tau^{-2}}{2}Q_j,\\
-\E_j\log q_j(\beta)
&=
\frac12\log(2\pi e\,v_{\beta j}),\\
-\E_j\log q_j(z_j)
&=
\frac12\log\left\{(2\pi e)^{J-1}|V_{zj}|\right\}.
\end{align}
These expressions define a numerically stable coordinate update by
computing \(\mathcal L_j\) for all \(j\) and applying a softmax after
subtracting \(\max_j\mathcal L_j\).

\section{Equivalent Scalar Form}

For implementation, one can avoid explicitly forming each \(S_j^{-1}\).
For any vector \(r\in\mathbb R^J\), define
\[
C_j(r)
=
(r_{-j}-\mu_jr_j)^\top S_j^{-1}(r_{-j}-\mu_jr_j).
\]
The block inverse identity implies
\[
C_j(r)=r^\top\Om r-r_j^2.
\]
When \(r=x\), set
\[
z_{\mathrm{denom},j}
=
N\overline{\beta^2}_j+\tau^{-2},
\qquad
z_{\mathrm{scale},j}
=
\frac{N\bar\beta_j}{z_{\mathrm{denom},j}}.
\]
After the \(q_j(z_j)\) update,
\begin{align}
Q_j
&=
\frac{J-1}{z_{\mathrm{denom},j}}
+z_{\mathrm{scale},j}^2C_j(x),\\
A_j&=1+Q_j,\\
B_j&=x_j+z_{\mathrm{scale},j}C_j(x).
\end{align}
This representation reduces the coordinate-specific calculations to
scalar operations once \(x^\top\Om x\) is available.

\section{Optional Empirical Bayes Update for \texorpdfstring{\(\nu_0\)}{nu0}}

Holding the variational distribution fixed, the \(\nu_0\)-dependent part
of the SER-relax ELBO is, up to constants,
\[
M(\nu_0)
=
\sum_{j=1}^J\pi_j
\left\{
\frac{J-1}{2}\log\tau^{-2}
-\frac{\tau^{-2}}{2}Q_j
\right\}.
\]
Because \(\sum_j\pi_j=1\), maximizing over \(\tau^2\) gives
\[
\hat\tau^2
=
\frac{\sum_{j=1}^J\pi_jQ_j}{J-1},
\qquad
\hat\nu_0=\frac{1}{\hat\tau^2}-3.
\]
In practice \(\nu_0\) should be constrained to be positive and bounded
away from numerically extreme values.

\end{document}
